
This dissertation developed a reproducible coarse-ROI occlusion audit workflow for low-resolution facial phenotyping within the Explainable Face2Gene Diagnosis project theme. On the supplied PD-DBS image-level split, the MLP and compact CNN captured strong label-associated signal. Low-level ROI summaries, similarity-threshold analyses, ROI-size effects, and YuNet sensitivity defined the main interpretation limits.


The central contribution is the conversion of model behaviour into structured ROI-level evidence that can be statistically compared, perturbed, and audited. The audit identified the data elements that would strengthen later clinical interpretation: patient identifiers, session metadata, acquisition records, clinical outcome scores, DBS and medication information, paired pre-/post structure, higher-resolution images, and external testing.


The project achieved its aim by producing a reproducible pipeline and structured audit outputs for the supplied pre-/post-DBS label task. The same ROI-evidence workflow can support later Face2Gene-style facial phenotyping or clinically anchored neurological facial-analysis tasks.


